一条递推只说了一件事：规模 \(n\) 的对象，由规模更小的同类对象拼出来。归并排序的 \(T(n)=2T(n/2)+cn\)、不含连续两个 1 的二进制串的 \(a_n=a_{n-1}+a_{n-2}\)、用面额 1、2、5 凑出 \(n\) 元的方案数，都是这个形状。规则写下来只占一行，你要的答案却是整条序列的走势：它涨得多快，第 \(n\) 项能不能一步算出，改动一个系数会带来什么。

这一章的主线是：解的形态由约化的方式决定。规模每次除以 \(b\)，展开出来是一棵越来越宽、深度只有 \(\log_b n\) 的树，总量由每层代价与层数这两个因子相乘得到。规模每次只减 1 或减 2，展开是一条不变宽的链，长期行为由一个代数方程的根说了算。而当拆分位置本身自由——规模 \(n\) 拆成 \(k\) 与 \(n-k\)，对所有 \(k\) 求和——前两种手段都够不着，得把整条序列打包成一个形式对象再动手。

\begin{center}
\begin{tikzpicture}[x=1cm,y=1cm,font=\footnotesize,>=stealth]
  \draw[black!55,fill=blue!10] (-2.4,0.55) rectangle (0.6,1.85);
  \node at (-0.9,1.45) {规模 \(n\) 的问题};
  \node at (-0.9,0.95) {约化为更小的同类};
  \draw[black!55,fill=blue!6] (2.6,2.55) rectangle (8.2,3.55);
  \node at (5.4,3.05) {除以常数 \(\Rightarrow\) 树：每层代价 \(\times\) 层数};
  \draw[black!55,fill=blue!6] (2.6,0.7) rectangle (8.2,1.7);
  \node at (5.4,1.2) {减去常数 \(\Rightarrow\) 链：特征根的幂};
  \draw[black!55,fill=blue!6] (2.6,-1.15) rectangle (8.2,-0.15);
  \node at (5.4,-0.65) {任意拆分 \(\Rightarrow\) 卷积：生成函数相乘};
  \draw[->,black!70] (0.7,1.35) -- (2.5,3.05);
  \draw[->,black!70] (0.7,1.2) -- (2.5,1.2);
  \draw[->,black!70] (0.7,1.05) -- (2.5,-0.65);
\end{tikzpicture}

\emph{同一句``由更小的同类拼成''，按三种方式约化，就要用三套不同的手法读出整体形态。三篇的顺序即按此排列。}
\end{center}

\begin{itemize}
\tightlist
\item \hyperref[sec:04-Discrete-Mathematics/02-Recurrences/01]{``递推建模与递归树''}：怎样从``最后一步有哪些互斥的可能''写出递推，以及分治开销为什么由每层代价与层数两项相乘决定；主定理的三种情形只是同一个等比级数的三种公比。
\item \hyperref[sec:04-Discrete-Mathematics/02-Recurrences/02]{``线性递推与特征根''}：常系数线性递推的解是特征根的幂的叠加。把递推写成向量迭代 \(v_{n+1}=Av_n\) 之后，这与矩阵幂由特征值决定就是同一件事，重根处冒出的 \(n\lambda^n\) 也和 Jordan 块对上了号。
\item \hyperref[sec:04-Discrete-Mathematics/02-Recurrences/03]{``生成函数把递推变成代数''}：把整条序列装进一个形式幂级数，下标平移变成乘 \(x\)，卷积变成乘法，递推塌缩成一个代数方程。换一个表示让问题变简单，这条思路后面还会以 Fourier 变换的样子再出现一次。
\end{itemize}

闭式并不是每次都值得追。很多场合下增长阶、稳定性或者一个能稳定迭代的递推已经够用，而一个漂亮的根式（Binet 公式就是例子）拿浮点数逐项算起来未必比原递推更可靠。动手之前先想清楚要的是哪一种答案：精确的第 \(n\) 项、增长的数量级，还是对结构的一句解释。
